\section{Co niezupełność oznacza dla codziennej matematyki}\label{sec:codzienna-matematyka}
\subsection*{„Nie da się udowodnić” nie znaczy „jest fałszywe”}\label{sec:16:1}
Dla konkretnego zdania $\varphi$ trzeba rozdzielić trzy komunikaty:
\begin{enumerate}[label=(\arabic*)]
\item $T\vdash\varphi$: istnieje dowód z ustalonych aksjomatów;
\item $T\nvdash\varphi$: takiego dowodu nie ma;
\item $\mathbb N\models\varphi$: zdanie jest prawdziwe o standardowych liczbach naturalnych.
\end{enumerate}
Drugi komunikat nie jest ani dowodem $\neg\varphi$, ani stwierdzeniem fałszywości $\varphi$ w $\mathbb N$. Jeżeli ani $\varphi$, ani $\neg\varphi$ nie ma dowodu w $T$, mówimy o niezależności od $T$, nie o braku wartości logicznej. W ustalonej strukturze i klasycznej semantyce zdanie nadal jest prawdziwe albo fałszywe.

Zdanie Gödla $G$ dobrze pokazuje różnicę: przy niesprzeczności $T$ jest prawdziwe w $\mathbb N$, lecz niedowodliwe w $T$ (Uwaga~\ref{res:9:3}). Ten wniosek uzyskujemy w metateorii, korzystając z założenia niesprzeczności. Nie otrzymaliśmy tym samym dowodu $G$ w $T$. Do niedowodliwości $\neg G$ wcześniejszy argument potrzebował silniejszej hipotezy; zdanie Rossera daje obustronną niezależność już przy zwykłej niesprzeczności. Nie wolno tych dwóch wersji zamieniać bez sprawdzenia założeń.

\subsection*{Dlaczego dodanie nowego aksjomatu nie kończy sprawy}\label{sec:16:2}
Załóżmy, że $T$ jest niesprzeczną, efektywnie aksjomatyzowalną teorią zawierającą $Q$, i weźmy jej niezależne zdanie Rossera $R$. Obie teorie $T+R$ oraz $T+\neg R$ są niesprzeczne. Gdyby pierwsza była sprzeczna, twierdzenie o dedukcji dawałoby $T\vdash\neg R$; gdyby druga była sprzeczna, dawałoby $T\vdash R$. Oba wnioski przeczą niezależności.

Można więc rozstrzygnąć ten jeden problem, dodając wybraną stronę jako aksjomat. Każde z tych rozszerzeń jest jednak nadal efektywnie aksjomatyzowalne i zawiera $Q$, zatem ma nowe zdanie niezależne. Nie chodzi o to, że to samo $R$ na zawsze pozostaje niedostępne: zmiana teorii zmienia również zdanie skonstruowane dla jej predykatu dowodliwości.

Także nieskończone dokładanie aksjomatów według jednej efektywnej procedury nie daje zupełnej teorii, jeśli wynikowa teoria pozostaje niesprzeczna i zawiera $Q$. Do całego efektywnie wyliczonego zbioru aksjomatów znów stosuje się twierdzenie Rossera. Można natomiast wziąć po prostu wszystkie prawdy $\operatorname{Th}(\mathbb N)$: otrzymamy teorię niesprzeczną i zupełną, ale nie efektywnie aksjomatyzowalną. Właśnie ten warunek wtedy tracimy.

Niesprzeczność obu rozszerzeń nie oznacza, że oba opisują prawdziwie standardowe $\mathbb N$: jedno zawiera $R$, a drugie jego negację. Co najwyżej jedno może mieć wszystkie aksjomaty prawdziwe w $\mathbb N$. Niesprzeczność jest słabszym wymaganiem niż prawdziwość wszystkich aksjomatów w zamierzonej strukturze.

\subsection*{Jak niezależność widać w modelach}\label{sec:16:3}
Lemat~\ref{res:14:7} daje model $T+R$ i model $T+\neg R$. Zatem te same wyjściowe aksjomaty $T$ dopuszczają światy różniące się odpowiedzią na pytanie $R$. Nie są to dwa sprzeczne opisy jednego ustalonego świata, lecz dwa modele aksjomatów, które nie rozstrzygnęły tej kwestii.

Tu widać, dlaczego zupełność logiki i niezupełność teorii nie przeczą sobie:
\begin{itemize}
\item \emph{zupełność logiki} mówi: jeśli zdanie zachodzi we wszystkich modelach $T$, istnieje jego dowód z $T$;
\item \emph{niezupełność teorii} mówi: dla pewnego zdania modele $T$ nie dają jednej wspólnej odpowiedzi --- istnieją modele obu stron.
\end{itemize}
W szczególności prawdziwość w jednym zamierzonym modelu $\mathcal M$ nie jest tym samym co prawdziwość we wszystkich modelach $T$.

\subsection*{Dlaczego nie istnieje uniwersalny przycisk „rozstrzygnij”}\label{sec:16:4}
Wniosek~\ref{res:9:6} wyklucza algorytm, który dla każdego zdania zawsze kończy pracę i poprawnie odpowiada, czy zdanie ma dowód w ustalonej niesprzecznej, efektywnie aksjomatyzowalnej $T\supseteq Q$. Jest to wynik o problemie decyzyjnym dowodliwości, nie samo inne sformułowanie niezupełności. W ogólności niezupełność teorii nie oznacza jeszcze nierozstrzygalności jej zbioru twierdzeń.

Można za to wypisywać kandydatów na dowody i sprawdzać ich poprawność. Jeśli poszukiwany dowód istnieje, takie przeszukiwanie kiedyś go znajdzie. Jeśli nie istnieje, samo przeszukiwanie może nigdy się nie zatrzymać. Równoległe poszukiwanie dowodu $\varphi$ i dowodu $\neg\varphi$ również nie pomaga, gdy zdanie jest niezależne: obie próby mogą trwać bez końca.

Nie oznacza to, że konkretne problemy są zwykle beznadziejne. Oznacza brak jednej zawsze kończącej się i zawsze poprawnej metody dla wszystkich wejść. Metody dla określonej klasy zadań, heurystyki oraz przeszukiwanie do ustalonej granicy pozostają możliwe. Brak rozwiązania przez długi czas sam w sobie nie dowodzi niezależności.

